%% kapa: exam=1 round=21 year=2010 period=2 form=A subject=accounting subject_ko=회계학 seq=25 number=25
%% source: Q-Net 감정평가사 기출문제 (cst00309) artlSeq=5024181 file=extracted/1008/5024181/2010년 제21회 감정평가사 1차시험 기출문제/2교시.hwp
%% conversion: kordoc 4.24.51 -> tools/gen1.py (auto); review=proofread-sonnet; flags=none
\begin{question}{25}
(주)감평은 20×9년 초 액면금액 ￦100,000의 3년 만기 전환사채를 액면금액으로 발행하였다. 전환권을 행사하지 않은 경우 전환사채의 만기일에 액면금액에 추가하여 상환할증금 ￦9,930을 지급한다. 전환사채의 표시이자율은 연 5\%이며 매년 말에 지급한다. (주)감평은 당사의 신용으로 만기 3년, 표시이자율 연 10\%, 액면금액으로 일반사채를 발행할 수 있다. 이 때 (주)감평이 자본으로 처리하는 전환권의 가치는 얼마인가? (단, 가장 근사치를 선택할 것)
\begin{center}
\small
\begin{tabular}{|c|c|c|}
\hline
\multirow{2}{*}{기간} & \multicolumn{2}{c|}{10\%} \\
\cline{2-3}
 & \makecell{1 회금액 ￦1의\\현재가치} & \makecell{기말연금 ￦1의\\현재가치} \\
\hline
1 & 0.90909 & 0.90909 \\
\hline
2 & 0.82645 & 1.73554 \\
\hline
3 & 0.75131 & 2.48685 \\
\hline
\end{tabular}
\end{center}
% [원문 그대로] 표의 '1 회금액'
\begin{choices}
  \item ￦4,974
  \item ￦9,739
  \item ￦12,434
  \item ￦16,232
  \item ￦25,971
\end{choices}
\end{question}
\answer{1}
